\documentclass[11pt,a4paper]{article}
\usepackage[T1]{fontenc}
\usepackage[utf8]{inputenc}
\usepackage{amsmath,amssymb,amsthm}
\usepackage[margin=1in]{geometry}
\usepackage[colorlinks=true,linkcolor=blue,urlcolor=blue]{hyperref}
\theoremstyle{plain}
\newtheorem{theorem}{Theorem}
\newtheorem{lemma}[theorem]{Lemma}
\newtheorem{proposition}[theorem]{Proposition}
\newtheorem{corollary}[theorem]{Corollary}
\theoremstyle{definition}
\newtheorem{remark}[theorem]{Remark}

\title{The empirical recurrence for OEIS A206088, proved by transfer matrix}
\author{Adrian Perez Fontelles\\ \small Independent researcher}
\date{2 September 2026}
\begin{document}
\maketitle

\begin{abstract}
OEIS A206088 counts arrays over $\{0,\dots,2\}$ under a condition on the perimeter cycle of
every $2\times3$ and $3\times2$ subblock, and carries an empirical recurrence of order $4$ contributed by
R.~H.~Hardin. It is true, and it is decidable rather than empirical. The condition is local to
a bounded window of consecutive rows, so the count is a walk count on $S=270$ states and is
$C$-finite; whether the conjectured recurrence annihilates it is then settled by a finite exact
computation, with the Cayley--Hamilton theorem bounding the work.
\end{abstract}

\noindent\small 2020 Mathematics Subject Classification. 05A15, 05B45, 15B36.\normalsize

\section{The conjecture}

OEIS A206088 is ``Number of (n+1) X 3 0..2 arrays with every 2 X 3 or 3 X 2 subblock having exactly two counterclockwise and two clockwise edge increases.''. It has offset $1$ and begins
\[
270, 780, 2016, 5952, 17976, 54432, 165936, 504912,\ \dots
\]
The entry states, as a conjecture contributed by its author and never marked settled:
\begin{quote}
Empirical: a(n) = a(n-1) + 6*a(n-2) + 4*a(n-3) - 10*a(n-4) for n>6.
\end{quote}
As of the ``Last modified'' line on the live entry (Jun 13 2018 11:23:42, revision 12) this is still
recorded as empirical, and nothing on the entry records it as proved.

\section{What is being counted}

Both a $3\times2$ and a $2\times3$ window have all six of their cells on the
perimeter, so each carries a $6$-cycle. For
$\begin{pmatrix}p&q\\r&s\\t&u\end{pmatrix}$ the clockwise cycle is
$p\to q\to s\to u\to t\to r\to p$; for
$\begin{pmatrix}p&q&t\\r&s&u\end{pmatrix}$ it is
$p\to q\to t\to u\to s\to r\to p$. An edge increase is a step of the cycle that goes up in
value, and counterclockwise counts the same cycle traversed the other way. The entry requires
of every window of either shape that it has exactly $2$ counterclockwise and $2$ clockwise edge increases.

The readings were pinned by direct enumeration before anything was built: over
$\{0,1,2\}$ the $3\times2$ arrays of clockwise count $2$ number $423$ and those of count
$3$ number $150$; over $\{0,\dots,3\}$ the $3\times2$ arrays of count $1$ number $480$;
and the $2\times3$ arrays over $\{0,1,2\}$ of clockwise count $2$ also number $423$. Those
are the first terms the corresponding entries publish.

\section{The count is a walk count}

A $3\times2$ window spans three consecutive array rows and a $2\times3$
window two, so take as state the PAIR of consecutive rows. The $2\times3$ windows lie inside a
single pair, so they decide which pairs are admissible at all; the $3\times2$ windows straddle
three rows, so they decide which pair may follow which. A step appends one array row, turning
$(r,s)$ into $(s,t)$ and testing every $3\times2$ window of $r,s,t$ at once. An $(n+1)$-row
array is a walk of $n-1$ steps and
\[
a(n)\;=\;\iota^{\!\top}M^{\,n-1}\tau,\qquad\iota=\tau=(1,\dots,1)^{\!\top},
\]
on the $S=270$ admissible pairs.

At $n=1$ the array has two rows, so no $3\times2$ window fits and only the
$2\times3$ windows apply; the entry's first term is
$270$, which the model reproduces along with every later term.
In particular $a$ is $C$-finite. The alignment of the walk index with the entry's own $n$ was
checked against its published terms rather than assumed.

\section{The proof}

\begin{lemma}\label{lem:crit}
Let $q(t)=t^{4}-\sum_i c_it^{4-i}$ be the characteristic polynomial of the
conjectured recurrence. The residual $a(n)-\sum_ic_ia(n-i)$ equals
$\iota^{\!\top}M^{\,j}q(M)\tau$ for the corresponding walk index $j$, so the recurrence
holds from some point on if and only if $\iota^{\!\top}M^{j}q(M)\tau=0$ for all large $j$,
and it is enough to examine $j<S$.
\end{lemma}

\begin{proof}
Factor $M^{j}$ out of $M^{j+4}-\sum_ic_iM^{j+4-i}$. For the bound, $M^{S}$
is an integer combination of $I,M,\dots,M^{S-1}$ by Cayley--Hamilton, so the numbers
$u_j=\iota^{\!\top}M^{j}q(M)\tau$ obey the monic recurrence given by the characteristic
polynomial of $M$; $S$ consecutive zeros force every later one, and the last nonzero $u_j$
pins the threshold exactly.
\end{proof}

\begin{theorem}
\[
a(n)\;=\;a(n-1) + 6\,a(n-2) + 4\,a(n-3) - 10\,a(n-4)
\]
for every $n>6$.
\end{theorem}

\begin{proof}
The vector $w=q(M)\tau$ was formed by $4$ matrix--vector products in exact integer
arithmetic and $\iota^{\!\top}M^{j}w$ evaluated for $j=0,1,\dots$, again exactly; those
integers vanish from the index corresponding to $n=6+1$ onwards and the run continued
until the working vector was identically zero, which settles every larger $j$ at once. The bound is exact: at $n=6$ the recurrence fails on the entry's own published terms, so no larger range holds.
\end{proof}

\section{Verification}

Three checks, each independent of the algebra above.

First, the model reproduces all $22$ terms the entry publishes, exactly, in integer
arithmetic. This is what ties the model to the entry: the digraph is built by reading the
entry's English, and a misreading gives different counts.

Second, the conjectured recurrence was evaluated directly on those published terms, with no
matrices involved, and holds wherever the proved range and the data overlap.

Third, the annihilation test of Lemma~\ref{lem:crit} was carried out in exact integer
arithmetic on the whole vector rather than on a sampled prefix.

\begin{thebibliography}{9}
\bibitem{oeis} The OEIS Foundation, \emph{The On-Line Encyclopedia of Integer
Sequences}, \url{https://oeis.org}, 2026. Sequence A206088.
\bibitem{stanley} R.~P.~Stanley, \emph{Enumerative Combinatorics}, Volume 1, 2nd ed.,
Cambridge University Press, 2011. (Transfer-matrix method, Section 4.7.)
\bibitem{bona} M.~B\'ona, \emph{Combinatorics of Permutations}, 2nd ed., CRC
Press, 2012.
\bibitem{horn} R.~A.~Horn and C.~R.~Johnson, \emph{Matrix Analysis}, 2nd ed., Cambridge
University Press, 2013.
\end{thebibliography}

\end{document}
